% Declaration sidecar for 10.1016_j.ejor.2020.03.034
% ASSISTED RESOLUTION — rung (c) of staged symbol resolution (sec:resolution).
% Written by corpusbuilder.assist, model deepseek-v4-flash, 2026-08-24.
% Non-deterministically sourced; pending human confirmation.
%
% Declaration sidecar for 10.1016_j.ejor.2020.03.034
% A multi-objective optimization-simulation approach for real time rescheduling in dense railway systems
%
% STUB — written by corpusbuilder.promote because no sidecar existed.
% Fill in the ? placeholders, delete the lines that do not apply, then save
% this file as corpus/declarations/10.1016_j.ejor.2020.03.034.tex and re-run promote.
%
% Every symbol below was read out of the accepted formulas. Where a line
% already carries a kind, the algebra or the reviewer settled it and the
% comment above it says which; check it, do not assume it. Where three lines
% appear, decide whether the symbol is an index family, a parameter or a
% variable and delete the two that are wrong. Bibliographic lines
% (meta/name/desc/prov) are generated from the dossier — do not add them here.
%
% index NAME  ordered=0|1 cyclic=0|1
% param NAME  shape=I,J|- kind=scalar|vector|matrix|big_m|tolerance
% var   NAME  shape=I,J|- domain=continuous|non_negative|integer|binary
%              role=primary|auxiliary|slack|indicator
% obj         sense=min|max name=IDENT combination=sum|weighted_sum|lexicographic
% con   NAME  kind=linear|big_m|ordering|headway|capacity|flow_balance|modulo|...

%@ obj sense=min name=objective combination=sum :: Minimizes recovery time, passenger quality of service, and number of proposed rescheduling decisions

% D: classified as index during review
%@ index D ordered=0 cyclic=0 :: Set of delayed events

% M: classified as index during review
%@ index M ordered=0 cyclic=0 :: Set of sufficiently large coefficients

% N: classified as index during review
%@ index N ordered=0 cyclic=0 :: Set of events

% P: classified as index during review
%@ index P ordered=0 cyclic=0 :: Set of planned times

% m: classified as index during review
%@ index m ordered=0 cyclic=0 :: Set of tolerance margins

% s: classified as index during review
%@ index s ordered=0 cyclic=0 :: Set of rescheduled times

% R: classified as parameter during review
%@ param R shape=- kind=scalar domain=- :: Recovery time

% T: classified as parameter during review
%@ param T shape=- kind=scalar domain=- :: Time

% time: classified as parameter during review
%@ param time shape=- kind=scalar domain=- :: Time component in quality of service

% In: classified as parameter during review
%@ param In shape=- kind=scalar domain=- :: In-vehicle time component

% QS: classified as parameter during review
%@ param QS shape=- kind=scalar domain=- :: Quality of service

% Waiting: classified as parameter during review
%@ param Waiting shape=- kind=scalar domain=- :: Waiting time component

% vehicle: classified as parameter during review
%@ param vehicle shape=- kind=scalar domain=- :: Vehicle time component

% unusable as an identifier, rename or drop: α
% Bound by a quantifier or big operator, so no declaration of their own: e, passengers
% (If one of these is really an index family, add a %@ index line for it.)

% Constraint rows (optional — they default to kind=linear):
%@ con eq_0001 kind=big_m domain=- indicator=- :: Defines when an event is delayed based on the tolerance margin and the big-M coefficient
%@ con eq_0002 kind=linear domain=- indicator=- :: Links recovery time to delayed events so that recovery time equals the latest delayed event time
%@ con eq_0003 kind=linear domain=- indicator=- :: Computes total passenger quality of service as weighted waiting time plus in-vehicle time
